% ===========================================================================
%  user/03-studio/00-chapter.tex
% ===========================================================================
\chapter{Studio}\label{cap:studio}

Studio is where things are built. You get there from \voce{Studio} in the
menu, or from the \pulsante{Open Studio} button on the home screen.

\figuraSchermata{1.0}{figure/schermate/studio.png}
  {Studio. At the top the three ways of working, below them the five families
   of artifacts, in the middle the drawing board.}
  {studio}

\section{The five families}\label{sez:studio-families}

The row of tabs below the title says \emph{what} you are building. They are
five different things, and it is worth knowing which one you need before you
start.

\begin{center}
\begin{tabular}{@{}L{28mm}L{110mm}@{}}
\toprule
\textbf{Tab} & \textbf{What it is} \\
\midrule
\rigaTenue
\voce{Strategies} & A strategy: it decides what to buy and sell, and when.
  It is the thing you build most often. \\
\voce{Approaches} & An approach: it skips the whole engine and produces the
  series of gains and losses directly. It is for people who already have a
  calculation of their own and just want to bring it into the system. \\
\rigaTenue
\voce{Weight mode} & A rule for sharing capital among several strategies.
  Fourteen come already installed. \\
\voce{Overlay} & A brake: it reduces exposure when the market behaves in a
  certain way. It does not earn, it protects. \\
\rigaTenue
\voce{Metrics} & A measure to add to the ones already shown on the result
  cards. \\
\bottomrule
\end{tabular}
\end{center}

\section{The three ways to build}\label{sez:studio-modes}

At the top right you choose how to work. The result is the same --- a code
file installed in the system --- and only the road changes.

\subsection{Blueprint: with blocks}

You drag blocks from the library on the left onto the board and connect
outputs to inputs. The library is divided into categories: prices, numbers,
external data such as earnings and fundamentals.

\begin{pratica}{The quickest way to see how it works}
Press \pulsante{Load an example}: an example appears, already wired (the
crossing of two moving averages with a stop). Take it apart and put it back
together: in ten minutes you understand the mechanism without writing a
line.
\end{pratica}

\subsection{Code: by writing}

The \pulsantesec{\{ \} Code} button switches to the text editor. It is the
same artifact, seen as Python code. You can switch between the two views.

\subsection{AI agent: in words}

The \voce{AI agent} tab opens a conversation: you describe the strategy in
words and the model writes it. It needs a model installed locally; without
one, the tab says so and the other two ways keep working.

\figuraSchermata{1.0}{figure/schermate/agent.png}
  {The assistant. You describe what you want, and it writes the artifact and
   installs it.}
  {agent}

\begin{attenzione}{Read what the assistant writes}
The generated code is installed in the system and runs on your data --- and
one day, perhaps, on your account. The assistant does not know whether a
strategy makes financial sense: it knows how to write code that does what it
was asked. The \voce{Code} tab is there for this too.
\end{attenzione}

\section{Genetics}\label{sez:studio-genetics}

The \voce{Genetics} tab opens an area of its own: instead of writing one
strategy, you let many of them be born and keep those that survive. You
choose the universe of instruments and let it run.

\figuraSchermata{1.0}{figure/schermate/genetic.png}
  {The genetics section. You declare what the strategies should evolve on,
   and the rest is waiting.}
  {genetic}

\begin{nota}{It takes a long time}
It is the slowest job in the application. It starts as a background process:
you can close the page, and the bell tells you when it is finished.
\end{nota}

\section{What happens when you press \enquote{Install}}\label{sez:studio-install}

The artifact is written to a folder of the system, and from then on it
appears in the lists: a strategy among those you can run, a weighting mode
among those you can choose in a project, an overlay among those you can
switch on.

Before writing, the system checks the code: it must not read or write files,
must not open connections and must not run commands. If it contains any of
these, the installation is refused and the reason is shown.
